\section{Introduction}
\label{sec:intro}

When an automated system denies someone a loan, a credit line, or a job, the
person is owed more than a score: they are owed an explanation they can
\emph{act} on. The dominant formalisation of such an explanation is the
\emph{counterfactual} of \citet{wachter2017counterfactual} --- the nearest input
$x'$ that the model would classify differently --- and a large literature now
generates these efficiently and diversely
\citep{mothilal2020dice,vanlooveren2021interpretable,verma2020counterfactual}.
Almost all of it optimises, and reports, a single property: \emph{validity}, that
$x'$ flips the model.

Validity is necessary. It is very far from sufficient. A counterfactual is a
piece of advice, and advice can flip a model while describing a world that could
never come about. We see all three of the following in the plans real methods
produce on standard data (\S\ref{sec:experiments}). A method tells an applicant
to acquire \$28{,}377 in capital gains --- valid, and useless. Another offers
additional capital \emph{losses} as a route to a higher income, exploiting a
correlation the model learned but no person can act on. A third flips the
decision by having the applicant work \emph{fewer} hours. Each is a valid
counterfactual under the standard definition. None is recourse.

These failures are not noise; they are two distinct, systematic gaps that
validity cannot see. The first is \emph{causal}: features cause one another, so
intervening on one moves others, and the record that actually results after a
person acts is not the record the method scored
\citep{karimi2021recourse,pearl2009causality}. Advice that ignores this is scored
in a world that acting on it destroys. The second is \emph{actionability}: a plan
may ask a person to change something immutable, protected, or only movable in one
direction \citep{ustun2019actionable,poyiadzi2020face}. Both gaps have been
studied in isolation; what has been missing is a single evaluation that holds a
plan to all of it at once, and a method built to pass.

\paragraph{The usable gap.} We separate three things a plan needs and treat them
as independent audits. A plan is \Valid{} if it flips the black box; \Causal{} if
it still flips the box once the structural causal model propagates the changes
into their consequences; and \Action{} if a person could actually carry it out.
We call a plan \Usable{} only when all three hold. The three are not nested: a
plan can be actionable yet causally invalid, or causally valid yet ask the
impossible. Measuring them separately is the point --- it turns ``does it flip the
model'' into ``is this advice you could follow and would it work,'' and it
reveals a gap the field has been reporting past. Figure~\ref{fig:usable_gap}
shows it: on our four datasets every method clears the validity bar, and the
\Usable{} rate nonetheless ranges from \mbox{13\%} to \mbox{100\%}.

\begin{figure}[t]
  \centering
  \includegraphics[width=0.86\linewidth]{results/results_usable_gap}
  \caption{The validity--usable gap. Faint bars: the share of rejected applicants
  a method gives a \emph{valid} plan (it flips the black box). Solid bars: the
  share it gives a \Usable{} plan (valid \emph{and} causally sound \emph{and}
  actionable). Every method clears validity; only \cac{} closes the gap on all
  four datasets. Numbers in Table~\ref{tab:main_results}.}
  \label{fig:usable_gap}
\end{figure}

\paragraph{Causal Actionable Counterfactuals.} We then close it. \cac{} searches
for an \emph{action} --- a small set of interventions $\dopart{X_S = v_S}$ --- and
scores each candidate not on the point it names but on the record its structural
causal model says the action would \emph{produce}, with every downstream feature
recomputed. Causal validity thereby becomes a property of the representation
rather than a penalty that can be traded away: \cac{} cannot propose a world its
own model disowns. Around that core it jointly minimises the effort of the
actions, their feasibility under an explicit actionability model, and the
plausibility of the resulting record, with a population search because the black
box is opaque and half the search space is a discrete subset choice.

\paragraph{Contributions.}
\begin{itemize}
  \item \textbf{A three-audit evaluation of recourse} (\S\ref{sec:desiderata})
    that separates validity, causal validity, and actionability, and scores
    methods on their conjunction, \Usable{} --- exposing a gap that
    validity-only evaluation hides (Figure~\ref{fig:usable_gap},
    Table~\ref{tab:main_results}).
  \item \textbf{The \cac{} method} (\S\ref{sec:method}): a gradient-free search
    over interventions, scored through a structural causal model, that is
    causally valid by construction and jointly optimises cost, feasibility, and
    plausibility.
  \item \textbf{A reproducible study on four public datasets}
    (\S\ref{sec:experiments}) showing the gap is real and large, that \cac{}
    closes it to 100\% usable recourse everywhere, that an ablation attributes
    the closing to the causal term, and that the cost of doing so is modest and
    honestly reported. Every dataset is public, every model runs locally, and the
    whole pipeline rebuilds from one command.
\end{itemize}
